\section{Grouping parties into two sides}
\label{sec:peps_party_coarsening}

Let $f:P\to Q$ be an arbitrary map of party labels. We group all parties
with the same image while retaining every register and every operation.
This is the owner identification used for the affected and exterior
systems in \cite[proof of Theorem~5.2]{OpenAI2026PolynomialPEPS},
\texttt{04-compression.tex}, lines 383--427.

% The two-side result below assumes the explicit internal-source condition.
% The chronological construction that frees corrected and crossing registers
% and retains untouched aggregate gates must still establish that condition.
% Neither that construction nor the full compression theorem is claimed here.

\begin{definition}[Changing register owners]
    \label{def:peps_owner_layout}
    \lean{TNLean.PEPS.PairEffect.Layout.mapOwner}
    \lean{TNLean.PEPS.PairEffect.Layout.mapOwner_nil}
    \lean{TNLean.PEPS.PairEffect.Layout.mapOwner_cons}
    \lean{TNLean.PEPS.PairEffect.Layout.mapOwner_append}
    \lean{TNLean.PEPS.PairEffect.Layout.mapOwnerIso}
    \lean{TNLean.PEPS.PairEffect.Layout.mapOwnerIso_nil_apply}
    \lean{TNLean.PEPS.PairEffect.Layout.mapOwnerIso_cons_tmul}
    \leanok
    \uses{def:peps_party_registers}
    If $\ell=((p_1,H_1),\ldots,(p_n,H_n))$ is a register list,
    let $\ell^f=((f(p_1),H_1),\ldots,(f(p_n),H_n))$.
    Empty lists and concatenation are preserved. The canonical isometry
    $R_{f,\ell}:\mathcal H_\ell\to\mathcal H_{\ell^f}$ is the
    tensor product of the identities on the original register spaces.
    On the empty list it is the identity of $\C$, and on an elementary
    tensor it satisfies
    \begin{align}
        R_{f,(p,H)\ell}(x\otimes y)
         & =x\otimes R_{f,\ell}y.
        \label{eq:peps_owner_head}
    \end{align}
    No injectivity or finiteness assumption on $f$, $P$, or $Q$ is imposed.
\end{definition}

\begin{theorem}[Compatibility with concatenation]
    \label{thm:peps_owner_concatenation}
    \lean{TNLean.PEPS.PairEffect.Layout.mapOwnerIso_append_tmul}
    \lean{TNLean.PEPS.PairEffect.Layout.mapOwnerIso_append}
    \leanok
    \uses{def:peps_owner_layout}
    Write $C_{\ell,m}:\mathcal H_{\ell m}\to
        \mathcal H_\ell\otimes\mathcal H_m$ for the canonical
    concatenation isometry. Under the identification
    $(\ell m)^f=\ell^f m^f$, for $x\in\mathcal H_\ell$
    and $y\in\mathcal H_m$ one has
    \begin{align}
        R_{f,\ell m}C_{\ell,m}^{-1}(x\otimes y)
         & =C_{\ell^f,m^f}^{-1}
        (R_{f,\ell}x\otimes R_{f,m}y).
        \label{eq:peps_owner_concatenation}
    \end{align}
    Equivalently,
    $R_{f,\ell m}C_{\ell,m}^{-1}
        =C_{\ell^f,m^f}^{-1}(R_{f,\ell}\otimes R_{f,m})$.
\end{theorem}

\begin{proof}\leanok
    Induct on $\ell$. For the empty list the assertion is scalar
    linearity. For a nonempty list, first consider a tensor with its
    leading register separated. Apply (\ref{eq:peps_owner_head})
    and the induction hypothesis to the remaining registers. Additivity
    extends the identity to arbitrary vectors. Equality on elementary
    tensors gives the operator identity.
\end{proof}

\begin{definition}[Sources remaining between distinct owners]
    \label{def:peps_owner_sources}
    \lean{TNLean.PEPS.PairEffect.PairSource.mapOwner}
    \lean{TNLean.PEPS.PairEffect.SourceInventory.mapOwner}
    \lean{TNLean.PEPS.PairEffect.SourceInventory.mapOwner_nil}
    \lean{TNLean.PEPS.PairEffect.SourceInventory.mapOwner_cons}
    \lean{TNLean.PEPS.PairEffect.SourceInventory.mapOwner_append}
    \lean{TNLean.PEPS.PairEffect.SourceInventory.mem_mapOwner}
    \lean{TNLean.PEPS.PairEffect.SourceInventory.map_partyPair_mapOwner}
    \lean{TNLean.PEPS.PairEffect.SourceInventory.mem_partyPair_mapOwner}
    \leanok
    \uses{def:peps_source_inventory}
    A source $s=(p,q;U,V;\eta)$ remains between distinct owners
    exactly when $f(p)\ne f(q)$. Its new record is then
    $(f(p),f(q);U,V;\eta)$, with the same spaces and vector.
    For an ordered source list $S$, let $S_f$ be the list of these
    surviving records in their original order. This operation preserves
    the empty list and concatenation.

    The unordered pair labels of $S_f$ are obtained by applying $f$
    to both endpoints of each original pair and removing the diagonal
    images. Their multiplicities are retained: different original pairs
    may acquire the same new label. Thus a pair $k$ occurs in $S_f$
    exactly when it is non-diagonal and is the image of some pair
    occurring in $S$.
\end{definition}

\begin{theorem}[Normalization and the absence of crossing sources]
    \label{thm:peps_owner_sources_internal}
    \lean{TNLean.PEPS.PairEffect.SourceInventory.IsNormalized.mapOwner}
    \lean{TNLean.PEPS.PairEffect.SourceInventory.mapOwner_eq_nil_iff}
    \leanok
    \uses{def:peps_owner_sources}
    If the vectors of $S$ are normalized, then those of $S_f$ are
    normalized. Moreover, $S_f$ is empty if and only if
    $f(p)=f(q)$ for the two endpoints of every source in $S$.
\end{theorem}

\begin{proof}\leanok
    Surviving records have unchanged vectors. A source is removed
    precisely when its two new owners agree, which proves the second
    assertion.
\end{proof}

\begin{definition}[Preparing an internal pair]
    \label{def:peps_internal_pair_preparation}
    \lean{TNLean.PEPS.PairEffect.Word.localPairSource}
    \leanok
    \uses{def:peps_party_layout}
    For $q\in Q$ and $\eta\in U\otimes V$, prepare two fresh
    registers $U,V$ at the single owner $q$ by the local map
    $c\mapsto\eta\otimes c$ from $\C$ to
    $U\otimes(V\otimes\C)$, with the canonical associativity
    identification. Tensor this map with the identity on all spectator
    registers.
\end{definition}

\begin{theorem}[The internal preparation operator]
    \label{thm:peps_internal_pair_preparation}
    \lean{TNLean.PEPS.PairEffect.Word.localPairSource_spec}
    \lean{TNLean.PEPS.PairEffect.Word.eval_localPairSource}
    \leanok
    \uses{def:peps_internal_pair_preparation}
    The operation in Definition~\ref{def:peps_internal_pair_preparation}
    maps a spectator vector $x$ to $\eta\otimes x$, under canonical
    associativity. If $\|\eta\|=1$, it is an allowed local operation
    and creates no pair source between distinct owners.
\end{theorem}

\begin{proof}\leanok
    The preparation map has norm at most $\|\eta\|$, and the
    associativity identification is isometric. Tensoring with the
    spectator identity preserves the contraction bound. For the operator
    identity, separate an elementary tensor in $\eta$ and apply the
    concatenation isometries; then extend by linearity.
\end{proof}

\begin{definition}[Changing the owners of a composition]
    \label{def:peps_owner_composition}
    \lean{TNLean.PEPS.PairEffect.Word.mapOwner}
    \leanok
    \uses{def:peps_owner_layout,def:peps_internal_pair_preparation}
    Let $F:\mathcal H_\ell\to\mathcal H_{\ell'}$ be an actual
    composition of local maps, pair preparations, exchanges, and
    operations beside untouched registers. Its composition $F^f$ on
    the new owners is defined in the same chronological order.
    A local map $A:\mathcal H_a\to\mathcal H_b$ becomes
    $R_{f,b}AR_{f,a}^{-1}$, acting at its new owner. Exchanges and
    untouched registers retain their original spaces. A pair source
    retains its vector and spaces; it remains a pair source when its
    new owners differ, and becomes the internal local preparation when
    they agree.
\end{definition}

\begin{theorem}[Exact preservation of an allowed composition]
    \label{thm:peps_owner_composition}
    \lean{TNLean.PEPS.PairEffect.Word.isAllowed_mapOwner}
    \lean{TNLean.PEPS.PairEffect.Word.sources_mapOwner}
    \lean{TNLean.PEPS.PairEffect.Word.eval_mapOwner}
    \lean{TNLean.PEPS.PairEffect.Word.eval_mapOwner_comp}
    \leanok
    \uses{def:peps_owner_composition,def:peps_owner_sources}
    The operators of $F$ and $F^f$ satisfy
    \begin{align}
        R_{f,\ell'}F & =F^fR_{f,\ell}.
        \label{eq:peps_owner_composition}
    \end{align}
    This identity does not require norm bounds. If $F$ is allowed,
    then $F^f$ is allowed. Its pair-source list is exactly $S_f$,
    where $S$ is the pair-source list of $F$.
\end{theorem}

\begin{proof}\leanok
    \uses{thm:peps_owner_concatenation,thm:peps_internal_pair_preparation}
    Induct on the chronological composition. For a local map,
    (\ref{eq:peps_owner_concatenation}) moves the owner identifications
    through its local and spectator factors; the conjugating isometries
    cancel. An internal source has the same preparation operator by
    Theorem~\ref{thm:peps_internal_pair_preparation}, while a surviving
    source has the same vector and endpoint spaces. Exchanges and
    operations beside untouched registers commute with the tensor
    product of identities. Composition preserves the equality.
    Isometric conjugation preserves each local contraction bound, and
    an internal normalized source is an allowed local preparation.
    The description of the source list follows by the same induction.
\end{proof}

\begin{theorem}[Two contractions when every source is internal]
    \label{thm:peps_internal_source_factorization}
    \lean{TNLean.PEPS.PairEffect.Word.exists_partition_of_sources_internal}
    \leanok
    % The statement below is an auxiliary result with its internal-source hypothesis explicit.
    \uses{def:peps_owner_composition,def:peps_party_partition}
    Let $f:P\to\{0,1\}$, and suppose $F$ is allowed. Assume that
    $f(p)=f(q)$ for the endpoints of every pair source occurring in
    $F$. Let $\mathcal H_{\ell,+}$ and $\mathcal H_{\ell,-}$ be
    the memories of the registers with owner images $1$ and $0$,
    respectively, in their original relative order. Let
    $J_{f,\ell}:\mathcal H_\ell\to
        \mathcal H_{\ell,+}\otimes\mathcal H_{\ell,-}$ first
    change the owner labels and then partition the two owners.
    There are two allowed compositions $A_+,A_-$ on the respective
    memories, each containing no pair sources and each having operator
    norm at most one, such that
    \begin{align}
        J_{f,\ell'}F
         & =(A_+\otimes A_-)J_{f,\ell}.
        \label{eq:peps_internal_source_factorization}
    \end{align}
    The two compositions are the actual restrictions of $F^f$ to
    the two owners.
\end{theorem}

\begin{proof}\leanok
    \uses{thm:peps_owner_sources_internal,thm:peps_owner_composition,
        thm:peps_party_binary_factorization,thm:peps_party_restriction}
    The internal-source condition makes $S_f$ empty. Hence $F^f$
    contains no pair preparations, while remaining an allowed
    composition. Apply Theorem~\ref{thm:peps_party_binary_factorization}
    to its two owners and compose the resulting identity with
    (\ref{eq:peps_owner_composition}). Its two restricted compositions
    are allowed and therefore have norm at most one.
\end{proof}
